% --- Archon physics formalization source begin ---
% archon:physics
% archon:covers PhyXMiniProblems/problem_phyx_mini_0909.lean
% archon:source-report reports/phyx_mini/problem_phyx_mini_0909.source.json

\chapter{Physics problem phyx\_mini\_0909}
\label{ch:PhyXMiniProblems_problem_phyx_mini_0909}

\paragraph{Problem source.}
\#\# Physical scenario

An electron gun creates a beam of electrons moving horizontally with a speed of \textbackslash{}( 3.3\textbackslash{}times10\textasciicircum{}\{7\}\textbackslash{},\textbackslash{}mathrm\{m/s\} \textbackslash{}). The electrons enter a \textbackslash{}( 2.0\textbackslash{},\textbackslash{}mathrm\{cm\} \textbackslash{})-long gap between two parallel electrodes where the electric field is \textbackslash{}( \textbackslash{}vec\{E\} = (5.0\textbackslash{}times10\textasciicircum{}\{4\}\textbackslash{},\textbackslash{}mathrm\{N/C\}, \textbackslash{}text\{down\}) \textbackslash{}).

\#\# Question

By what angle, is the electron beam deflected by these electrodes?

\#\# Answer choices

- A: A: 3\textasciicircum{}\textbackslash{}circ

- B: B: 9\textasciicircum{}\textbackslash{}circ

- C: C: 9.1\textasciicircum{}\textbackslash{}circ

- D: D: 20\textasciicircum{}\textbackslash{}circ

\#\# Figure caption (auxiliary; use the image as primary evidence)

The image illustrates a scenario involving charged deflection plates and an electric field. 

1. **Deflection Plates**: 
   - Two parallel plates are shown, one above the other. The top plate has positive charges depicted by plus signs, and the bottom plate has a dashed line.

2. **Electric Field (\textbackslash{}(\textbackslash{}vec\{E\}\textbackslash{}))**: 
   - The electric field is indicated between the plates, pointing downwards with magenta arrows. It is labeled as \textbackslash{}(\textbackslash{}vec\{E\} = (5.0 \textbackslash{}times 10\textasciicircum{}4 \textbackslash{}, \textbackslash{}text\{N/C, down\})\textbackslash{}).

3. **Particle Path**:
   - A particle enters the electric field horizontally from the left with initial velocity \textbackslash{}(\textbackslash{}vec\{v\}\_0\textbackslash{}) (indicated by a green arrow). It follows a curved path upward as it exits the field, ending with velocity \textbackslash{}(\textbackslash{}vec\{v\}\_1\textbackslash{}) (blue arrow).
   - The angle of deflection is indicated as \textbackslash{}(\textbackslash{}theta\textbackslash{}).

4. **Distance**:
   - The length of the plates is labeled \textbackslash{}(L = 2.0 \textbackslash{}, \textbackslash{}text\{cm\}\textbackslash{}).

5. **Text**:
   - The lower part of the image includes the text "Deflection plates."

The diagram likely represents the deflection of a charged particle due to an electric field within the plates, demonstrating the change in vector direction caused by the field.

\#\# Dataset metadata

Electrostatics; Physical Model Grounding Reasoning, Multi-Formula Reasoning

\paragraph{Recorded answer/context.}
C: C: 9.1\textasciicircum{}\textbackslash{}circ

\paragraph{Figure/image path.}
/root/proposal\_for\_physic/science-mango/phyx\_mini\_run/phyx\_data/test\_image/909.png

\paragraph{Formalization target.}
create a compiling Lean file with sorry bodies at `PhyXMiniProblems/problem\_phyx\_mini\_0909.lean`. The Lean declarations must preserve the physical quantities, dimensions or dimensional roles, figure labels, governing-law hypotheses, and final relation expressed by this problem.
Use Mathlib/Physlib names found through LeanExplore where available. If a physics API is missing, introduce faithful local abstractions rather than scalar placeholder aliases.

\begin{theorem}[Physics formalization target]
\label{thm:physics:phyx_mini_0909:target}
\lean{PhyXMiniProblems.ProblemPhyXMini0909.problem_phyx_mini_0909}
\uses{def:physics:phyx-mini-0909:phyxminiproblems-problemphyxmini0909-electrondeflectionsetup, def:physics:phyx-mini-0909:phyxminiproblems-problemphyxmini0909-matcheselectrondeflectionscenario, def:physics:phyx-mini-0909:phyxminiproblems-problemphyxmini0909-matchessuppliedelectrondeflectionfigure, def:physics:phyx-mini-0909:phyxminiproblems-problemphyxmini0909-matchesgivenelectrondeflectionreadouts, def:physics:phyx-mini-0909:phyxminiproblems-problemphyxmini0909-usesstandardelectronreferencedata, def:physics:phyx-mini-0909:phyxminiproblems-problemphyxmini0909-hasphysicalelectrondeflectionparameters, def:physics:phyx-mini-0909:phyxminiproblems-problemphyxmini0909-satisfiesuniformelectricdeflectiondynamics, def:physics:phyx-mini-0909:phyxminiproblems-problemphyxmini0909-deflectionangleradians, def:physics:phyx-mini-0909:phyxminiproblems-problemphyxmini0909-deflectionangledegrees, def:physics:phyx-mini-0909:phyxminiproblems-problemphyxmini0909-predicteddeflectionangleradians, def:physics:phyx-mini-0909:phyxminiproblems-problemphyxmini0909-recordeddatasetanswer, def:physics:phyx-mini-0909:phyxminiproblems-problemphyxmini0909-isuniqueclosestdisplayedangle}
The assigned autoformalize agent should translate this physics problem into Lean declarations in the covered file, with theorem and lemma proof bodies written as `by sorry`.
\end{theorem}
\begin{proof}
This is an autoformalization task, not a proof task. Produce statements that can later be proved by the physics prover without weakening the physical model.
\end{proof}

% --- Archon declaration topology begin ---
\section{Declaration topology}
\begin{definition}[velocity Dimension]
  \label{def:physics:phyx-mini-0909:phyxminiproblems-problemphyxmini0909-velocitydimension}
  \lean{PhyXMiniProblems.ProblemPhyXMini0909.velocityDimension}
  The physical dimension L T⁻¹ of velocity
\end{definition}
\begin{proof}
  This is the declared model definition of velocity Dimension.
\end{proof}
\begin{definition}[acceleration Dimension]
  \label{def:physics:phyx-mini-0909:phyxminiproblems-problemphyxmini0909-accelerationdimension}
  \lean{PhyXMiniProblems.ProblemPhyXMini0909.accelerationDimension}
  The physical dimension L T⁻² of acceleration
\end{definition}
\begin{proof}
  This is the declared model definition of acceleration Dimension.
\end{proof}
\begin{definition}[force Dimension]
  \label{def:physics:phyx-mini-0909:phyxminiproblems-problemphyxmini0909-forcedimension}
  \lean{PhyXMiniProblems.ProblemPhyXMini0909.forceDimension}
  The physical dimension M L T⁻² of force
\end{definition}
\begin{proof}
  This is the declared model definition of force Dimension.
\end{proof}
\begin{definition}[electric Field Dimension]
  \label{def:physics:phyx-mini-0909:phyxminiproblems-problemphyxmini0909-electricfielddimension}
  \lean{PhyXMiniProblems.ProblemPhyXMini0909.electricFieldDimension}
  The physical dimension M L T⁻² C⁻¹ of electric field
\end{definition}
\begin{proof}
  This is the declared model definition of electric Field Dimension.
\end{proof}
\begin{definition}[Mass Quantity]
  \label{def:physics:phyx-mini-0909:phyxminiproblems-problemphyxmini0909-massquantity}
  \lean{PhyXMiniProblems.ProblemPhyXMini0909.MassQuantity}
  A positive, unit-independent physical mass
\end{definition}
\begin{proof}
  This is the declared model definition of Mass Quantity.
\end{proof}
\begin{definition}[Signed Charge Quantity]
  \label{def:physics:phyx-mini-0909:phyxminiproblems-problemphyxmini0909-signedchargequantity}
  \lean{PhyXMiniProblems.ProblemPhyXMini0909.SignedChargeQuantity}
  A signed, unit-independent electric charge
\end{definition}
\begin{proof}
  This is the declared model definition of Signed Charge Quantity.
\end{proof}
\begin{definition}[Length Quantity]
  \label{def:physics:phyx-mini-0909:phyxminiproblems-problemphyxmini0909-lengthquantity}
  \lean{PhyXMiniProblems.ProblemPhyXMini0909.LengthQuantity}
  A nonnegative, unit-independent physical length
\end{definition}
\begin{proof}
  This is the declared model definition of Length Quantity.
\end{proof}
\begin{definition}[Time Quantity]
  \label{def:physics:phyx-mini-0909:phyxminiproblems-problemphyxmini0909-timequantity}
  \lean{PhyXMiniProblems.ProblemPhyXMini0909.TimeQuantity}
  A nonnegative, unit-independent elapsed time
\end{definition}
\begin{proof}
  This is the declared model definition of Time Quantity.
\end{proof}
\begin{definition}[Planar Velocity Quantity]
  \label{def:physics:phyx-mini-0909:phyxminiproblems-problemphyxmini0909-planarvelocityquantity}
  \lean{PhyXMiniProblems.ProblemPhyXMini0909.PlanarVelocityQuantity}
  \uses{def:physics:phyx-mini-0909:phyxminiproblems-problemphyxmini0909-velocitydimension}
  A unit-independent velocity vector in the plane of the figure
\end{definition}
\begin{proof}
  This is the declared model definition of Planar Velocity Quantity.
\end{proof}
\begin{definition}[Planar Acceleration Quantity]
  \label{def:physics:phyx-mini-0909:phyxminiproblems-problemphyxmini0909-planaraccelerationquantity}
  \lean{PhyXMiniProblems.ProblemPhyXMini0909.PlanarAccelerationQuantity}
  \uses{def:physics:phyx-mini-0909:phyxminiproblems-problemphyxmini0909-accelerationdimension}
  A unit-independent acceleration vector in the plane of the figure
\end{definition}
\begin{proof}
  This is the declared model definition of Planar Acceleration Quantity.
\end{proof}
\begin{definition}[Planar Force Quantity]
  \label{def:physics:phyx-mini-0909:phyxminiproblems-problemphyxmini0909-planarforcequantity}
  \lean{PhyXMiniProblems.ProblemPhyXMini0909.PlanarForceQuantity}
  \uses{def:physics:phyx-mini-0909:phyxminiproblems-problemphyxmini0909-forcedimension}
  A unit-independent force vector in the plane of the figure
\end{definition}
\begin{proof}
  This is the declared model definition of Planar Force Quantity.
\end{proof}
\begin{definition}[Planar Electric Field Quantity]
  \label{def:physics:phyx-mini-0909:phyxminiproblems-problemphyxmini0909-planarelectricfieldquantity}
  \lean{PhyXMiniProblems.ProblemPhyXMini0909.PlanarElectricFieldQuantity}
  \uses{def:physics:phyx-mini-0909:phyxminiproblems-problemphyxmini0909-electricfielddimension}
  A unit-independent electric-field vector in the plane of the figure
\end{definition}
\begin{proof}
  This is the declared model definition of Planar Electric Field Quantity.
\end{proof}
\begin{definition}[x Axis]
  \label{def:physics:phyx-mini-0909:phyxminiproblems-problemphyxmini0909-xaxis}
  \lean{PhyXMiniProblems.ProblemPhyXMini0909.xAxis}
  Coordinate 0, horizontal and positive toward the right in the figure
\end{definition}
\begin{proof}
  This is the declared model definition of x Axis.
\end{proof}
\begin{definition}[y Axis]
  \label{def:physics:phyx-mini-0909:phyxminiproblems-problemphyxmini0909-yaxis}
  \lean{PhyXMiniProblems.ProblemPhyXMini0909.yAxis}
  Coordinate 1, vertical and positive upward in the figure
\end{definition}
\begin{proof}
  This is the declared model definition of y Axis.
\end{proof}
\begin{definition}[mass In Kilograms]
  \label{def:physics:phyx-mini-0909:phyxminiproblems-problemphyxmini0909-massinkilograms}
  \lean{PhyXMiniProblems.ProblemPhyXMini0909.massInKilograms}
  \uses{def:physics:phyx-mini-0909:phyxminiproblems-problemphyxmini0909-massquantity}
  Read a mass in coherent-SI kilograms
\end{definition}
\begin{proof}
  This is the declared model definition of mass In Kilograms.
\end{proof}
\begin{definition}[charge In Coulombs]
  \label{def:physics:phyx-mini-0909:phyxminiproblems-problemphyxmini0909-chargeincoulombs}
  \lean{PhyXMiniProblems.ProblemPhyXMini0909.chargeInCoulombs}
  \uses{def:physics:phyx-mini-0909:phyxminiproblems-problemphyxmini0909-signedchargequantity}
  Read a signed charge in coherent-SI coulombs
\end{definition}
\begin{proof}
  This is the declared model definition of charge In Coulombs.
\end{proof}
\begin{definition}[length In Meters]
  \label{def:physics:phyx-mini-0909:phyxminiproblems-problemphyxmini0909-lengthinmeters}
  \lean{PhyXMiniProblems.ProblemPhyXMini0909.lengthInMeters}
  \uses{def:physics:phyx-mini-0909:phyxminiproblems-problemphyxmini0909-lengthquantity}
  Read a length in coherent-SI metres
\end{definition}
\begin{proof}
  This is the declared model definition of length In Meters.
\end{proof}
\begin{definition}[length In Centimeters]
  \label{def:physics:phyx-mini-0909:phyxminiproblems-problemphyxmini0909-lengthincentimeters}
  \lean{PhyXMiniProblems.ProblemPhyXMini0909.lengthInCentimeters}
  \uses{def:physics:phyx-mini-0909:phyxminiproblems-problemphyxmini0909-lengthquantity, def:physics:phyx-mini-0909:phyxminiproblems-problemphyxmini0909-lengthinmeters}
  Read a length in centimetres, as used by the plate-length label
\end{definition}
\begin{proof}
  This is the declared model definition of length In Centimeters.
\end{proof}
\begin{definition}[time In Seconds]
  \label{def:physics:phyx-mini-0909:phyxminiproblems-problemphyxmini0909-timeinseconds}
  \lean{PhyXMiniProblems.ProblemPhyXMini0909.timeInSeconds}
  \uses{def:physics:phyx-mini-0909:phyxminiproblems-problemphyxmini0909-timequantity}
  Read an elapsed time in coherent-SI seconds
\end{definition}
\begin{proof}
  This is the declared model definition of time In Seconds.
\end{proof}
\begin{definition}[velocity Vector In Meters Per Second]
  \label{def:physics:phyx-mini-0909:phyxminiproblems-problemphyxmini0909-velocityvectorinmeterspersecond}
  \lean{PhyXMiniProblems.ProblemPhyXMini0909.velocityVectorInMetersPerSecond}
  \uses{def:physics:phyx-mini-0909:phyxminiproblems-problemphyxmini0909-planarvelocityquantity}
  Read a velocity vector in coherent-SI metres per second
\end{definition}
\begin{proof}
  This is the declared model definition of velocity Vector In Meters Per Second.
\end{proof}
\begin{definition}[acceleration Vector In Meters Per Second Squared]
  \label{def:physics:phyx-mini-0909:phyxminiproblems-problemphyxmini0909-accelerationvectorinmeterspersecondsquared}
  \lean{PhyXMiniProblems.ProblemPhyXMini0909.accelerationVectorInMetersPerSecondSquared}
  \uses{def:physics:phyx-mini-0909:phyxminiproblems-problemphyxmini0909-planaraccelerationquantity}
  Read an acceleration vector in coherent-SI metres per second squared
\end{definition}
\begin{proof}
  This is the declared model definition of acceleration Vector In Meters Per Second Squared.
\end{proof}
\begin{definition}[force Vector In Newtons]
  \label{def:physics:phyx-mini-0909:phyxminiproblems-problemphyxmini0909-forcevectorinnewtons}
  \lean{PhyXMiniProblems.ProblemPhyXMini0909.forceVectorInNewtons}
  \uses{def:physics:phyx-mini-0909:phyxminiproblems-problemphyxmini0909-planarforcequantity}
  Read a force vector in coherent-SI newtons
\end{definition}
\begin{proof}
  This is the declared model definition of force Vector In Newtons.
\end{proof}
\begin{definition}[electric Field Vector In Newtons Per Coulomb]
  \label{def:physics:phyx-mini-0909:phyxminiproblems-problemphyxmini0909-electricfieldvectorinnewtonspercoulomb}
  \lean{PhyXMiniProblems.ProblemPhyXMini0909.electricFieldVectorInNewtonsPerCoulomb}
  \uses{def:physics:phyx-mini-0909:phyxminiproblems-problemphyxmini0909-planarelectricfieldquantity}
  Read an electric-field vector in coherent-SI newtons per coulomb
\end{definition}
\begin{proof}
  This is the declared model definition of electric Field Vector In Newtons Per Coulomb.
\end{proof}
\begin{definition}[Particle Species]
  \label{def:physics:phyx-mini-0909:phyxminiproblems-problemphyxmini0909-particlespecies}
  \lean{PhyXMiniProblems.ProblemPhyXMini0909.ParticleSpecies}
  The charged particle species named in the prose
\end{definition}
\begin{proof}
  This is the declared model definition of Particle Species.
\end{proof}
\begin{definition}[Deflection Plate]
  \label{def:physics:phyx-mini-0909:phyxminiproblems-problemphyxmini0909-deflectionplate}
  \lean{PhyXMiniProblems.ProblemPhyXMini0909.DeflectionPlate}
  The two parallel deflection plates shown in the image
\end{definition}
\begin{proof}
  This is the declared model definition of Deflection Plate.
\end{proof}
\begin{definition}[Plate Charge Sign]
  \label{def:physics:phyx-mini-0909:phyxminiproblems-problemphyxmini0909-platechargesign}
  \lean{PhyXMiniProblems.ProblemPhyXMini0909.PlateChargeSign}
  Charge signs visibly printed on the two electrodes
\end{definition}
\begin{proof}
  This is the declared model definition of Plate Charge Sign.
\end{proof}
\begin{definition}[Beam State]
  \label{def:physics:phyx-mini-0909:phyxminiproblems-problemphyxmini0909-beamstate}
  \lean{PhyXMiniProblems.ProblemPhyXMini0909.BeamState}
  Entry and exit states carrying the velocity labels v₀ and v₁
\end{definition}
\begin{proof}
  This is the declared model definition of Beam State.
\end{proof}
\begin{definition}[Arrow Direction]
  \label{def:physics:phyx-mini-0909:phyxminiproblems-problemphyxmini0909-arrowdirection}
  \lean{PhyXMiniProblems.ProblemPhyXMini0909.ArrowDirection}
  Qualitative arrow directions appearing in the primary image
\end{definition}
\begin{proof}
  This is the declared model definition of Arrow Direction.
\end{proof}
\begin{definition}[Figure Object]
  \label{def:physics:phyx-mini-0909:phyxminiproblems-problemphyxmini0909-figureobject}
  \lean{PhyXMiniProblems.ProblemPhyXMini0909.FigureObject}
  Objects whose presence is semantically relevant in image 909.png
\end{definition}
\begin{proof}
  This is the declared model definition of Figure Object.
\end{proof}
\begin{definition}[Figure Label]
  \label{def:physics:phyx-mini-0909:phyxminiproblems-problemphyxmini0909-figurelabel}
  \lean{PhyXMiniProblems.ProblemPhyXMini0909.FigureLabel}
  Literal labels printed in image 909.png
\end{definition}
\begin{proof}
  This is the declared model definition of Figure Label.
\end{proof}
\begin{definition}[Electric Field Model]
  \label{def:physics:phyx-mini-0909:phyxminiproblems-problemphyxmini0909-electricfieldmodel}
  \lean{PhyXMiniProblems.ProblemPhyXMini0909.ElectricFieldModel}
  The idealized field model between the electrodes
\end{definition}
\begin{proof}
  This is the declared model definition of Electric Field Model.
\end{proof}
\begin{definition}[Electron Deflection Figure]
  \label{def:physics:phyx-mini-0909:phyxminiproblems-problemphyxmini0909-electrondeflectionfigure}
  \lean{PhyXMiniProblems.ProblemPhyXMini0909.ElectronDeflectionFigure}
  \uses{def:physics:phyx-mini-0909:phyxminiproblems-problemphyxmini0909-deflectionplate, def:physics:phyx-mini-0909:phyxminiproblems-problemphyxmini0909-platechargesign, def:physics:phyx-mini-0909:phyxminiproblems-problemphyxmini0909-arrowdirection, def:physics:phyx-mini-0909:phyxminiproblems-problemphyxmini0909-figureobject, def:physics:phyx-mini-0909:phyxminiproblems-problemphyxmini0909-figurelabel}
  Literal and qualitative content transcribed from the primary raster. Its numerical fields are printed labels, not definitions of the physical velocity or the requested angle
\end{definition}
\begin{proof}
  This is the declared model definition of Electron Deflection Figure.
\end{proof}
\begin{definition}[Electron Deflection Setup]
  \label{def:physics:phyx-mini-0909:phyxminiproblems-problemphyxmini0909-electrondeflectionsetup}
  \lean{PhyXMiniProblems.ProblemPhyXMini0909.ElectronDeflectionSetup}
  \uses{def:physics:phyx-mini-0909:phyxminiproblems-problemphyxmini0909-massquantity, def:physics:phyx-mini-0909:phyxminiproblems-problemphyxmini0909-signedchargequantity, def:physics:phyx-mini-0909:phyxminiproblems-problemphyxmini0909-lengthquantity, def:physics:phyx-mini-0909:phyxminiproblems-problemphyxmini0909-timequantity, def:physics:phyx-mini-0909:phyxminiproblems-problemphyxmini0909-planarvelocityquantity, def:physics:phyx-mini-0909:phyxminiproblems-problemphyxmini0909-planaraccelerationquantity, def:physics:phyx-mini-0909:phyxminiproblems-problemphyxmini0909-planarforcequantity, def:physics:phyx-mini-0909:phyxminiproblems-problemphyxmini0909-planarelectricfieldquantity, def:physics:phyx-mini-0909:phyxminiproblems-problemphyxmini0909-particlespecies, def:physics:phyx-mini-0909:phyxminiproblems-problemphyxmini0909-beamstate, def:physics:phyx-mini-0909:phyxminiproblems-problemphyxmini0909-electricfieldmodel, def:physics:phyx-mini-0909:phyxminiproblems-problemphyxmini0909-electrondeflectionfigure}
  Independent physical quantities of the beam passage. In particular, the exit velocity is not defined from the desired angle, and the transit time, acceleration, and force are not defined from an answer choice
\end{definition}
\begin{proof}
  This is the declared model definition of Electron Deflection Setup.
\end{proof}
\begin{definition}[Matches Electron Deflection Scenario]
  \label{def:physics:phyx-mini-0909:phyxminiproblems-problemphyxmini0909-matcheselectrondeflectionscenario}
  \lean{PhyXMiniProblems.ProblemPhyXMini0909.MatchesElectronDeflectionScenario}
  \uses{def:physics:phyx-mini-0909:phyxminiproblems-problemphyxmini0909-electrondeflectionsetup}
  The prose model: an electron traverses a uniform field between parallel plates
\end{definition}
\begin{proof}
  This is the declared model definition of Matches Electron Deflection Scenario.
\end{proof}
\begin{definition}[Matches Supplied Electron Deflection Figure]
  \label{def:physics:phyx-mini-0909:phyxminiproblems-problemphyxmini0909-matchessuppliedelectrondeflectionfigure}
  \lean{PhyXMiniProblems.ProblemPhyXMini0909.MatchesSuppliedElectronDeflectionFigure}
  \uses{def:physics:phyx-mini-0909:phyxminiproblems-problemphyxmini0909-yaxis, def:physics:phyx-mini-0909:phyxminiproblems-problemphyxmini0909-lengthincentimeters, def:physics:phyx-mini-0909:phyxminiproblems-problemphyxmini0909-electricfieldvectorinnewtonspercoulomb, def:physics:phyx-mini-0909:phyxminiproblems-problemphyxmini0909-electrondeflectionsetup}
  Primary-image evidence, including both plate polarities, the downward field, the two velocity arrows, the upward-curving path, θ, and the two printed numerical labels. It contains no value for θ
\end{definition}
\begin{proof}
  This is the declared model definition of Matches Supplied Electron Deflection Figure.
\end{proof}
\begin{definition}[Matches Given Electron Deflection Readouts]
  \label{def:physics:phyx-mini-0909:phyxminiproblems-problemphyxmini0909-matchesgivenelectrondeflectionreadouts}
  \lean{PhyXMiniProblems.ProblemPhyXMini0909.MatchesGivenElectronDeflectionReadouts}
  \uses{def:physics:phyx-mini-0909:phyxminiproblems-problemphyxmini0909-xaxis, def:physics:phyx-mini-0909:phyxminiproblems-problemphyxmini0909-yaxis, def:physics:phyx-mini-0909:phyxminiproblems-problemphyxmini0909-lengthincentimeters, def:physics:phyx-mini-0909:phyxminiproblems-problemphyxmini0909-velocityvectorinmeterspersecond, def:physics:phyx-mini-0909:phyxminiproblems-problemphyxmini0909-electricfieldvectorinnewtonspercoulomb, def:physics:phyx-mini-0909:phyxminiproblems-problemphyxmini0909-electrondeflectionsetup}
  Numerical physical readouts supplied in the prose and figure. The coordinate convention makes downward field components negative and the horizontal entry velocity positive
\end{definition}
\begin{proof}
  This is the declared model definition of Matches Given Electron Deflection Readouts.
\end{proof}
\begin{definition}[Uses Standard Electron Reference Data]
  \label{def:physics:phyx-mini-0909:phyxminiproblems-problemphyxmini0909-usesstandardelectronreferencedata}
  \lean{PhyXMiniProblems.ProblemPhyXMini0909.UsesStandardElectronReferenceData}
  \uses{def:physics:phyx-mini-0909:phyxminiproblems-problemphyxmini0909-massinkilograms, def:physics:phyx-mini-0909:phyxminiproblems-problemphyxmini0909-chargeincoulombs, def:physics:phyx-mini-0909:phyxminiproblems-problemphyxmini0909-electrondeflectionsetup}
  Standard reference values implicit in the word “electron.” They are physical input data, not a premise about the requested angle
\end{definition}
\begin{proof}
  This is the declared model definition of Uses Standard Electron Reference Data.
\end{proof}
\begin{definition}[Has Physical Electron Deflection Parameters]
  \label{def:physics:phyx-mini-0909:phyxminiproblems-problemphyxmini0909-hasphysicalelectrondeflectionparameters}
  \lean{PhyXMiniProblems.ProblemPhyXMini0909.HasPhysicalElectronDeflectionParameters}
  \uses{def:physics:phyx-mini-0909:phyxminiproblems-problemphyxmini0909-xaxis, def:physics:phyx-mini-0909:phyxminiproblems-problemphyxmini0909-massinkilograms, def:physics:phyx-mini-0909:phyxminiproblems-problemphyxmini0909-chargeincoulombs, def:physics:phyx-mini-0909:phyxminiproblems-problemphyxmini0909-lengthinmeters, def:physics:phyx-mini-0909:phyxminiproblems-problemphyxmini0909-timeinseconds, def:physics:phyx-mini-0909:phyxminiproblems-problemphyxmini0909-velocityvectorinmeterspersecond, def:physics:phyx-mini-0909:phyxminiproblems-problemphyxmini0909-electrondeflectionsetup}
  Positivity and nondegeneracy conditions for the intended physical branch
\end{definition}
\begin{proof}
  This is the declared model definition of Has Physical Electron Deflection Parameters.
\end{proof}
\begin{definition}[Satisfies Uniform Electric Deflection Dynamics]
  \label{def:physics:phyx-mini-0909:phyxminiproblems-problemphyxmini0909-satisfiesuniformelectricdeflectiondynamics}
  \lean{PhyXMiniProblems.ProblemPhyXMini0909.SatisfiesUniformElectricDeflectionDynamics}
  \uses{def:physics:phyx-mini-0909:phyxminiproblems-problemphyxmini0909-xaxis, def:physics:phyx-mini-0909:phyxminiproblems-problemphyxmini0909-massinkilograms, def:physics:phyx-mini-0909:phyxminiproblems-problemphyxmini0909-chargeincoulombs, def:physics:phyx-mini-0909:phyxminiproblems-problemphyxmini0909-lengthinmeters, def:physics:phyx-mini-0909:phyxminiproblems-problemphyxmini0909-timeinseconds, def:physics:phyx-mini-0909:phyxminiproblems-problemphyxmini0909-velocityvectorinmeterspersecond, def:physics:phyx-mini-0909:phyxminiproblems-problemphyxmini0909-accelerationvectorinmeterspersecondsquared, def:physics:phyx-mini-0909:phyxminiproblems-problemphyxmini0909-forcevectorinnewtons, def:physics:phyx-mini-0909:phyxminiproblems-problemphyxmini0909-electricfieldvectorinnewtonspercoulomb, def:physics:phyx-mini-0909:phyxminiproblems-problemphyxmini0909-electrondeflectionsetup}
  The elementary nonrelativistic model in coherent SI units. The first two clauses are F = qE and F = ma; the third is the constant-acceleration velocity update; the fourth says that constant horizontal motion traverses the plate length during the field-interaction time. None mentions an angle or answer choice
\end{definition}
\begin{proof}
  This is the declared model definition of Satisfies Uniform Electric Deflection Dynamics.
\end{proof}
\begin{definition}[deflection Angle Radians]
  \label{def:physics:phyx-mini-0909:phyxminiproblems-problemphyxmini0909-deflectionangleradians}
  \lean{PhyXMiniProblems.ProblemPhyXMini0909.deflectionAngleRadians}
  \uses{def:physics:phyx-mini-0909:phyxminiproblems-problemphyxmini0909-velocityvectorinmeterspersecond, def:physics:phyx-mini-0909:phyxminiproblems-problemphyxmini0909-electrondeflectionsetup}
  The physical deflection angle is the undirected Euclidean angle, in radians, between the independent entry and exit velocity vectors. This geometric definition does not insert a numerical answer
\end{definition}
\begin{proof}
  This is the declared model definition of deflection Angle Radians.
\end{proof}
\begin{definition}[radians To Degrees]
  \label{def:physics:phyx-mini-0909:phyxminiproblems-problemphyxmini0909-radianstodegrees}
  \lean{PhyXMiniProblems.ProblemPhyXMini0909.radiansToDegrees}
  Convert a real radian readout to degrees
\end{definition}
\begin{proof}
  This is the declared model definition of radians To Degrees.
\end{proof}
\begin{definition}[deflection Angle Degrees]
  \label{def:physics:phyx-mini-0909:phyxminiproblems-problemphyxmini0909-deflectionangledegrees}
  \lean{PhyXMiniProblems.ProblemPhyXMini0909.deflectionAngleDegrees}
  \uses{def:physics:phyx-mini-0909:phyxminiproblems-problemphyxmini0909-electrondeflectionsetup, def:physics:phyx-mini-0909:phyxminiproblems-problemphyxmini0909-deflectionangleradians, def:physics:phyx-mini-0909:phyxminiproblems-problemphyxmini0909-radianstodegrees}
  The deflection angle of the beam, read in degrees
\end{definition}
\begin{proof}
  This is the declared model definition of deflection Angle Degrees.
\end{proof}
\begin{definition}[predicted Deflection Angle Radians]
  \label{def:physics:phyx-mini-0909:phyxminiproblems-problemphyxmini0909-predicteddeflectionangleradians}
  \lean{PhyXMiniProblems.ProblemPhyXMini0909.predictedDeflectionAngleRadians}
  \uses{def:physics:phyx-mini-0909:phyxminiproblems-problemphyxmini0909-xaxis, def:physics:phyx-mini-0909:phyxminiproblems-problemphyxmini0909-yaxis, def:physics:phyx-mini-0909:phyxminiproblems-problemphyxmini0909-massinkilograms, def:physics:phyx-mini-0909:phyxminiproblems-problemphyxmini0909-chargeincoulombs, def:physics:phyx-mini-0909:phyxminiproblems-problemphyxmini0909-lengthinmeters, def:physics:phyx-mini-0909:phyxminiproblems-problemphyxmini0909-velocityvectorinmeterspersecond, def:physics:phyx-mini-0909:phyxminiproblems-problemphyxmini0909-electricfieldvectorinnewtonspercoulomb, def:physics:phyx-mini-0909:phyxminiproblems-problemphyxmini0909-electrondeflectionsetup}
  The arctangent prediction obtained from F = qE, F = ma, the traversal time L / v₀, and tan θ = v\_y / v\_x. It depends only on independent physical inputs and is not definitionally equal to the geometric angle above
\end{definition}
\begin{proof}
  This is the declared model definition of predicted Deflection Angle Radians.
\end{proof}
\begin{definition}[Answer Choice]
  \label{def:physics:phyx-mini-0909:phyxminiproblems-problemphyxmini0909-answerchoice}
  \lean{PhyXMiniProblems.ProblemPhyXMini0909.AnswerChoice}
  Labels of the four displayed answer choices
\end{definition}
\begin{proof}
  This is the declared model definition of Answer Choice.
\end{proof}
\begin{definition}[displayed Angle Degrees]
  \label{def:physics:phyx-mini-0909:phyxminiproblems-problemphyxmini0909-answerchoice-displayedangledegrees}
  \lean{PhyXMiniProblems.ProblemPhyXMini0909.AnswerChoice.displayedAngleDegrees}
  \uses{def:physics:phyx-mini-0909:phyxminiproblems-problemphyxmini0909-answerchoice}
  Angle in degrees printed beside each displayed answer label
\end{definition}
\begin{proof}
  This is the declared model definition of displayed Angle Degrees.
\end{proof}
\begin{definition}[recorded Dataset Answer]
  \label{def:physics:phyx-mini-0909:phyxminiproblems-problemphyxmini0909-recordeddatasetanswer}
  \lean{PhyXMiniProblems.ProblemPhyXMini0909.recordedDatasetAnswer}
  \uses{def:physics:phyx-mini-0909:phyxminiproblems-problemphyxmini0909-answerchoice}
  The answer label recorded in the supplied dataset metadata
\end{definition}
\begin{proof}
  This is the declared model definition of recorded Dataset Answer.
\end{proof}
\begin{definition}[Is Closest Displayed Angle]
  \label{def:physics:phyx-mini-0909:phyxminiproblems-problemphyxmini0909-isclosestdisplayedangle}
  \lean{PhyXMiniProblems.ProblemPhyXMini0909.IsClosestDisplayedAngle}
  \uses{def:physics:phyx-mini-0909:phyxminiproblems-problemphyxmini0909-answerchoice}
  A displayed choice is at least as close as every alternative
\end{definition}
\begin{proof}
  This is the declared model definition of Is Closest Displayed Angle.
\end{proof}
\begin{definition}[Is Unique Closest Displayed Angle]
  \label{def:physics:phyx-mini-0909:phyxminiproblems-problemphyxmini0909-isuniqueclosestdisplayedangle}
  \lean{PhyXMiniProblems.ProblemPhyXMini0909.IsUniqueClosestDisplayedAngle}
  \uses{def:physics:phyx-mini-0909:phyxminiproblems-problemphyxmini0909-answerchoice, def:physics:phyx-mini-0909:phyxminiproblems-problemphyxmini0909-isclosestdisplayedangle}
  A displayed choice is the unique closest one
\end{definition}
\begin{proof}
  This is the declared model definition of Is Unique Closest Displayed Angle.
\end{proof}
\begin{lemma}[exit velocity components]
  \label{lem:physics:phyx-mini-0909:phyxminiproblems-problemphyxmini0909-exit-velocity-components}
  \lean{PhyXMiniProblems.ProblemPhyXMini0909.exit_velocity_components}
  \uses{def:physics:phyx-mini-0909:phyxminiproblems-problemphyxmini0909-xaxis, def:physics:phyx-mini-0909:phyxminiproblems-problemphyxmini0909-yaxis, def:physics:phyx-mini-0909:phyxminiproblems-problemphyxmini0909-massinkilograms, def:physics:phyx-mini-0909:phyxminiproblems-problemphyxmini0909-chargeincoulombs, def:physics:phyx-mini-0909:phyxminiproblems-problemphyxmini0909-timeinseconds, def:physics:phyx-mini-0909:phyxminiproblems-problemphyxmini0909-velocityvectorinmeterspersecond, def:physics:phyx-mini-0909:phyxminiproblems-problemphyxmini0909-electricfieldvectorinnewtonspercoulomb, def:physics:phyx-mini-0909:phyxminiproblems-problemphyxmini0909-electrondeflectionsetup, def:physics:phyx-mini-0909:phyxminiproblems-problemphyxmini0909-hasphysicalelectrondeflectionparameters, def:physics:phyx-mini-0909:phyxminiproblems-problemphyxmini0909-satisfiesuniformelectricdeflectiondynamics}
  The force and kinematics laws first determine the exit-velocity components: the horizontal component is unchanged, while the vertical component is generated by the upward electric acceleration of the negatively charged electron
\end{lemma}
\begin{proof}
  The conclusion follows from the governing relations and figure readouts named in the statement of exit velocity components.
\end{proof}
\begin{lemma}[deflection angle eq arctangent prediction]
  \label{lem:physics:phyx-mini-0909:phyxminiproblems-problemphyxmini0909-deflection-angle-eq-arctangent-prediction}
  \lean{PhyXMiniProblems.ProblemPhyXMini0909.deflection_angle_eq_arctangent_prediction}
  \uses{def:physics:phyx-mini-0909:phyxminiproblems-problemphyxmini0909-electrondeflectionsetup, def:physics:phyx-mini-0909:phyxminiproblems-problemphyxmini0909-matchesgivenelectrondeflectionreadouts, def:physics:phyx-mini-0909:phyxminiproblems-problemphyxmini0909-usesstandardelectronreferencedata, def:physics:phyx-mini-0909:phyxminiproblems-problemphyxmini0909-hasphysicalelectrondeflectionparameters, def:physics:phyx-mini-0909:phyxminiproblems-problemphyxmini0909-satisfiesuniformelectricdeflectiondynamics, def:physics:phyx-mini-0909:phyxminiproblems-problemphyxmini0909-deflectionangleradians, def:physics:phyx-mini-0909:phyxminiproblems-problemphyxmini0909-predicteddeflectionangleradians}
  Combining the exit-velocity relation with the horizontal transit time gives the arctangent formula for the geometric deflection angle
\end{lemma}
\begin{proof}
  The conclusion follows from the governing relations and figure readouts named in the statement of deflection angle eq arctangent prediction.
\end{proof}
% --- Archon declaration topology end ---
% --- Archon physics formalization source end ---
